\section{Discussion}
From the results presented, it can be deduced that there may well be huge potential to exploit chemical saturation effects to reduce the non-\ce{CO_2} warming from aviation. Intersecting plumes in high-density airspace regions can lead to a diminishing influence of both \ce{NO_x}-based effects and persistent contrail growth. In this study, high-density hotspots were first investigated through a global heatmap analysis (figure \ref{WorldPlot}). This plot of ``mean airspace density'', combined with the peak aircraft counts per hour, derived in figure \ref{Hist}, exemplify the temporal inhomogeneity of air traffic distribution, typically seen at aircraft cruising altitudes. Even for the highest density grid cells, the maximum average aircraft density is just above 0.6 aircraft per grid volume. Comparing this to a maximum flux of 31 aircraft traversing through the EU hotspot, at its peak hour shows that traffic patterns can be highly nonlinear and unpredictable.

In the plume plots of figure \ref{PlumeUS} and \ref{PlumeEU}, it is evident that, despite the EU hotspot having a higher number of aircraft pairings over the 3-hour period, the traffic in the US hotspot seems to be following a much more organised track structure. This has a large effect on the ability to achieve plume overlap, as can be seen in the results in table \ref{tab}. Even though EU has 329 more instances of aircraft pairings recorded, the vast majority of its data points did not record an SPE occurrence (81.4\%). US on the other hand, achieved a 42\% SPE success rate, despite having a smaller aircraft count. As a result, this suggests that the US hotspot could have serve as a better candidate for controlled plume overlap, and that there are many other factors in play that may determine the effectiveness of this method (e.g. airspace complexity, temporal distribution, airspace organisation etc.).